\paragraph{树的行话。}\label{sub:tree-speak}
涉及多个\trPointedDisc{}的\trGluing{}运算，自然要借\trTree{}来描述。我们在此把所有相关的术语汇编
起来，其中有些是为后文备用的。所谓\textbf{$d$ \trLeafedTree}（$d$-leafed tree），$d>0$，指的是一棵%
\trProper\trEmbedding{}的平面\trTree{}$T\subset\RR^2$，它有 $d+1$ 条半无限\trEdge：其中一条被单独挑出，称为\trRoot，
其余的则是\trLeaf。半无限\trEdge{}也称\trExteriorEdge，其余的则称\trInteriorEdge。我们以 $\mathrm{Ve}(T)$ 表示%
\trVertex{}的集合，以 $\mathrm{Ed}^{\mathrm{int}}(T)$ 表示\trInteriorEdge{}的集合。回顾一下，图（graph）中的一个\trFlag{}是由
一个\trVertex{}与一条相邻的\trEdge{}组成的对 $f=(v,e)$。我们以 $\mathrm{Fl}^{\mathrm{int}}(T)$ 表示\trInteriorFlag{}的
集合（即其中\trEdge{}为\trInteriorEdge{}的那些旗）。
我们将按从\trRoot{}流向诸\trLeaf{}的“向上”方向为 $T$ \trOrientation，见\cref{fig:leafed-tree}。若 $e$ 的\trOrientation{}指向
背离 $v$ 的一侧，则称\trFlag{}$f=(v,e)$ 为\textbf{正}（positive），否则称其为\textbf{负}（negative）。每条%
\trInteriorEdge{}$e$ 都对应两个旗 $f^{\pm}(e)$，一正一负。类似地，每条\trExteriorEdge{}只有归属于它的一个旗，
这个旗为正还是为负，取决于该\trEdge{}是\trLeaf{}还是\trRoot。偶尔我们会发现，把与给定\trVertex{}$v$ 相邻的
诸旗编号为 $f_0(v),\ldots,f_{|v|-1}(v)$ 颇为有用，其中 $|v|$ 是\trValency。这一编号总是从唯一的负旗 $f_0(v)$
开始，随后按给定的平面\trEmbedding{}\trAnticlockwise{}地延续。
\begin{figure}[htbp]
\centering
\begin{tikzpicture}[line cap=round,line join=round]
\node at (-1.62,0.30) {根};
\draw[->] (-1.70,0) -- (0,0);
\draw[->] (0,0) -- (3.09,0.85);
\draw[->] (0,0) -- (1.79,-1.21);
\draw[->] (1.79,-1.21) -- (3.08,-0.73);
\draw[->] (1.79,-1.21) -- (3.08,-1.99);
\draw (3.45,0.92) .. controls (3.74,0.90) and (3.50,0.02) .. (3.66,-0.60);
\draw (3.45,-2.06) .. controls (3.74,-2.04) and (3.50,-1.24) .. (3.66,-0.60);
\node at (4.28,-0.58) {叶};
\end{tikzpicture}
\caption{}\label{fig:leafed-tree}
\end{figure}
\emph{$d$ \trLeafedTree}称为\textbf{稳定}（stable）的，若所有\trVertex{}$v$ 都满足 $|v|\ge 3$；称为\textbf{半稳定}
（semistable）的，若 $|v|\ge 2$。
